%& -shell-escape
% !TeX root = thesis-abstract.tex
% !TeX encoding = UTF-8
% !TeX program = XeLaTeX

\documentclass{.local/lib/classes/ndvr-super-article-standalone}

\begin{document}
    
    Выбранная сигнатура представляет собой временной ряд.
    Существует большое количество алгоритмов для индексации
    и поиска временных рядов.
    В работе предлагается модифицированный вариант 
    алгоритма \textit{$\iSAX{2.0}$}\footnote{
        A. Camerra, T. Palpanas, J. Shieh, and E. Keogh. 
        iSax2.0: Indexing and mining one billion time series. 
        In ICDM, pages 58–67, 2010.
    }.
    Алгоритм $\iSAX{2.0}$ — улучшенная модификация алгоритма
    $\iSAX$\footnote{
        J. Shieh and E. Keogh. 
        iSax: Indexing and mining terabyte sized time series. 
        In SIGKDD, pages 623–631, 2008.
    },
    который основан на алгоритме 
    $\SAX$\footnote{
        J. Lin, E. Keogh, S. Lonardi, and B. Chiu, 
        A Symbolic Representation of Time Series, 
        with Implications for Streaming Algorithms, 
        Workshop on Research Issues in Data Mining and
        Knowledge Discovery, 
        the 8th ACM SIGMOD. San Diego, CA, 2003. 
    }.

    Алгоритмы 
    $\SAX$\footnote{
        $\SAX$ —  \textbf{S}ymbolic \textbf{A}ggregate  
        Appro\textbf{X}imation. 
        Перевод с английского: 
        «Символьная агрегатная аппроксимация».
    }
    и
    $\iSAX$\footnote{
        $\iSAX$ — \textbf{I}ndexable \textbf{S}ymbolic 
        \textbf{A}ggregate  Appro\textbf{X}imation.
        Перевод с английского 
        «Индексируемая символьная агрегатная аппроксимация».
    } 
    используют кусочно-агрегатную аппроксимацию 
    временного ряда. Каждый временной ряд представляется
    как последовательность символов.
    При переходе к символьной форме временной ряд 
    претерпевает дискретизацию и квантование.
    Границы квантования определяются в~соответствии 
    со стандартным нормальным распределением.

\end{document}
